\chapter*{Abstract}

Biomethane upgrading plants rely on a vacuum pump as a critical component of the vacuum pressure swing
adsorption (VPSA) process, where an unplanned failure interrupts production for the entire plant. Condition
monitoring of such assets is usually decomposed into three related tasks, anomaly detection, fault
detection and predictive maintenance, which the literature typically addresses in isolation, each with a
dedicated model. The few multi-task approaches proposed to date combine at most two of the three, and none
has been applied to biomethane upgrading, where no prior work addresses the vacuum pump specifically.

This dissertation develops and evaluates a multi-task deep learning framework that performs the three tasks
jointly, from a single shared representation of the sensor data. One month of real operational data was
obtained from a VPSA plant in Lithuania, comprising fourteen process variables sampled at one-minute
intervals. Since no vacuum pump failure has ever been recorded across the operator's monitored plants, the
data contains only normal operation. A cleaning pipeline was therefore developed to partition the record
into eleven independent production sessions, and a fault injection framework was built on top of it,
generating four fault scenarios derived from the manufacturer's troubleshooting guide at three severity
levels, two scenarios in which faults overlap in time, and a false positive scenario. Each injected
deviation is calibrated to the manufacturer's warning and trip thresholds and propagated to the sensors
physically correlated with the affected one, yielding a labelled dataset of 145{,}826 observations that
supports all three tasks simultaneously.

Five models were compared across eleven window sizes, from a single observation to 1{,}024: Random Forest
and XGBoost as traditional baselines, three single-task LSTM models, and two multi-task LSTM architectures
differing in whether the task-specific heads branch independently off the shared backbone or are chained,
each conditioned on the predictions of the preceding ones.

The recurrent models outperformed the traditional baselines on all three tasks, most markedly on predictive
maintenance, where the baselines failed to improve on a constant predictor at most window sizes while the
deep learning models reached a coefficient of determination of 0.678. The single-task models obtained the
best anomaly detection F1-score (0.900) and the best predictive maintenance $R^2$ (0.678), while the
cascaded multi-task architecture obtained the best fault detection mean F1-score (0.720), with margins of
at most 0.023 in every case. A single multi-task model therefore performs comparably to three independently
trained ones while covering all three tasks in one forward pass, which favours the multi-task formulation
for a deployed monitoring system.

\vspace*{10mm}\noindent
\textbf{Keywords}: Multi-task learning, Anomaly detection, Fault detection, Predictive maintenance,
Long Short-Term Memory, Fault injection, Biomethane upgrading, Vacuum pump

\chapter*{Resumo}

As centrais de purificação de biometano dependem de uma bomba de vácuo como componente crítico do processo
de adsorção por variação de pressão em vácuo (VPSA), no qual uma falha não planeada interrompe a produção
de toda a instalação. A monitorização da condição destes equipamentos é habitualmente decomposta em três
tarefas relacionadas, deteção de anomalias, deteção de falhas e manutenção preditiva, que a literatura
aborda tipicamente de forma isolada, com um modelo dedicado a cada uma. As poucas abordagens multitarefa
propostas até à data combinam no máximo duas das três, e nenhuma foi aplicada à purificação de biometano,
onde não existe trabalho anterior dedicado especificamente à bomba de vácuo.

Esta dissertação desenvolve e avalia uma abordagem de aprendizagem profunda multitarefa que executa as três
tarefas em conjunto, a partir de uma única representação partilhada dos dados dos sensores. Foi obtido um
mês de dados operacionais reais de uma central VPSA na Lituânia, com catorze variáveis de processo
amostradas ao minuto. Uma vez que nunca foi registada qualquer avaria da bomba de vácuo nas centrais
monitorizadas pelo operador, os dados contêm apenas operação normal. Foi por isso desenvolvido um processo
de limpeza que particiona o registo em onze sessões de produção independentes, sobre o qual foi construído
um sistema de injeção de falhas que gera quatro cenários derivados do guia de resolução de problemas do
fabricante, em três níveis de severidade, dois cenários com falhas sobrepostas no tempo e um cenário de
falso positivo. Cada desvio injetado é calibrado pelos limiares de aviso e de disparo do fabricante e
propagado aos sensores fisicamente correlacionados com o afetado, resultando num conjunto de dados
etiquetado com 145{,}826 observações que suporta as três tarefas em simultâneo.

Foram comparados cinco modelos ao longo de onze dimensões de janela temporal, de uma única observação até
1{,}024: Random Forest e XGBoost como referências tradicionais, três modelos LSTM de tarefa única e duas
arquiteturas LSTM multitarefa, que diferem entre si consoante as cabeças específicas de cada tarefa derivam
independentemente da representação partilhada ou são encadeadas, cada uma condicionada pelas previsões das
anteriores.

Os modelos recorrentes superaram as referências tradicionais nas três tarefas, de forma mais acentuada na
manutenção preditiva, onde as referências não conseguiram superar um preditor constante na maioria das
dimensões de janela, enquanto os modelos de aprendizagem profunda alcançaram um coeficiente de determinação
de 0,678. Os modelos de tarefa única obtiveram o melhor F1 na deteção de anomalias (0,900) e o melhor $R^2$
na manutenção preditiva (0,678), ao passo que a arquitetura multitarefa encadeada obteve o melhor F1 médio
na deteção de falhas (0,720), com margens não superiores a 0,023 em todos os casos. Um único modelo
multitarefa apresenta assim um desempenho comparável ao de três modelos treinados de forma independente,
cobrindo as três tarefas numa única passagem pela rede, o que favorece a formulação multitarefa para um
sistema de monitorização em produção.

\vspace*{10mm}\noindent
\textbf{Palavras-chave}: Aprendizagem multitarefa, Deteção de anomalias, Deteção de falhas,
Manutenção preditiva, Long Short-Term Memory, Injeção de falhas, Purificação de biometano, Bomba de vácuo